\documentclass[a4paper,zihao=-4]{ctexart}

\usepackage{geometry}
\usepackage{fancyhdr}
\usepackage{indentfirst}
\usepackage{setspace}
\usepackage{array}
\usepackage{etoolbox}
\usepackage{styles/joufonts}
\usepackage{styles/jouheadings}

\geometry{
    a4paper,
    top=2.5cm,
    bottom=2.5cm,
    left=2.5cm,
    right=2.5cm,
    headheight=0.8cm,
    headsep=0.5cm,
    footskip=1cm
}

\linespread{1.25}
\setlength{\parindent}{2em}
\setlength{\parskip}{0pt}
\newlength{\samplecontentwidth}
\setlength{\samplecontentwidth}{14.63cm}
\newlength{\samplexoffset}
\setlength{\samplexoffset}{0.68cm}

\newcommand{\samplebodypagestyle}{%
    \fancyhf{}
    \fancyhead[C]{%
        \makebox[\textwidth][c]{%
            \kaishu\zihao{5}%
            \begin{tabular*}{\samplecontentwidth}{@{}l@{\extracolsep{\fill}}c r@{}}
                \JOUHeadingBodyHeaderLeftInline & 理工农医类 & \JOUHeadingBodyHeaderRightPlaceholder
            \end{tabular*}%
        }%
    }
    \fancyfoot[C]{\zihao{5}\thepage}
    \renewcommand{\headrulewidth}{0.4pt}
    \renewcommand{\footrulewidth}{0pt}
}

\pagestyle{fancy}
\samplebodypagestyle

\newcommand{\samplebodytext}[1]{%
    {\songti\zihao{-4}\noindent #1\par}%
}

\begin{document}
\setcounter{page}{42}
\thispagestyle{fancy}

\vspace*{0.03cm}
\noindent\hspace*{\samplexoffset}\begin{minipage}[t]{\samplecontentwidth}
\JOUChapterHeadingLine{1}{绪论（一级标题）}
\vspace{0.16cm}
\samplebodytext{\hspace*{2em}×××××××××××××正文×××××××××××××××××××××××××××××…………}
\vspace{0.16cm}

\JOUSectionHeadingLine{1.1}{××××××（二级标题）}
\vspace{0.10cm}
\samplebodytext{\hspace*{2em}×××××××××××××正文×××××××××××××××××××××××…………}
\vspace{0.10cm}

\JOUSubsectionHeadingLine{1.1.1}{××××（三级标题）}
\vspace{0.10cm}
\samplebodytext{\hspace*{2em}××××××××××××正文××××××××××××××××××××…………}
\vspace{0.06cm}
\samplebodytext{\hspace*{2em}…………}
\vspace{1.08cm}

\JOUChapterHeadingLine{2}{×××××（一级标题）}
\vspace{0.18cm}
\samplebodytext{\hspace*{2em}×××××××××××××正文×××××××××××××××××××××××××××××…………}
\vspace{0.42cm}

{\songti\zihao{5}\noindent\mbox{*注：}\par}
\vspace{0.12cm}
{\songti\zihao{5}
\noindent\hspace*{2em}1. 本页为正文式样，本注释不是正文的部分，只是本式样的说明解释；\par
\noindent\hspace*{2em}2．标题编号方法应采用分级阿拉伯数字编号方法，第一级为“1”“2”“3”等，\par
\noindent 第二级为“2.1”“2.2”“2.3”等，第三级为“2.2.1”“2.2.2”“2.2.3”等，但分级阿拉伯数字\par
\noindent 的编号一般不超过四级，两级之间用下角圆点隔开，每一级的末尾不加标点，编\par
\noindent 号与标题内容间空一个汉字空格，标题应顶格；四级以下时可以先使用\par
\noindent {\fontspec{Arial Unicode MS}⑴⑵⑶……后①②③……}编制小标题，格式同正文；\par
\noindent\hspace*{2em}3. 一级标题为小三号黑体，缩放、间距、位置标准，无首行缩进，无左右\par
\noindent 缩进，行间距 1.25 倍多倍行距，段前、段后各 0.5 行间距；\par
\noindent\hspace*{2em}4. 二级标题为四号黑体，缩放、间距、位置标准，无首行缩进，无左右缩\par
\noindent 进，行间距 1.25 倍多倍行距，段前、段后无间距；\par
\noindent\hspace*{2em}5. 三级以下标题为小四号黑体，缩放、间距、位置标准，无首行缩进，无\par
\noindent 左右缩进，行间距 1.25 倍多倍行距，段前、段后无间距；\par
\noindent\hspace*{2em}6. 正文在标题下另起段不空行，为小四号，中文用宋体，英文用 Times New\par
\noindent Roman，缩放、间距、位置标准，无左右缩进，首行缩进 2 字符（两个汉字），\par
\noindent 无悬挂式缩进，段前、段后间距无，行间距为 1.25 倍多倍行距；\par
\noindent\hspace*{2em}7. 正文中表格与插图的中文一律用五号楷体\_GB2312，英文用五号 Times\par
\noindent New Roman；表格用三线表；\par
\noindent\hspace*{2em}8. 页眉字号为五号，中文用楷体\_GB2312，英文用 Times New Roman；\par
\noindent\hspace*{2em}9. 正文每一章开始，应另起一页。分页时请用插入分页符方式。\par
}
\end{minipage}

\end{document}
